%--------------------------------------------------------------------------------------------------
% 
\chapter*{Abstract}
\pdfbookmark[0]{Abstract}{Abstract}
%--------------------------------------------------------------------------------------------------
The recent development of large language models (LLMs) has significantly advanced the field of Natural Language Processing (NLP) and taken a leading role in the development of Artificial Intelligence.
We can see that trends are shifting not only towards extremely large models with multiple modalities and massive numbers of parameters but also towards smaller models that can be deployed outside large datacenter clusters.
However, smaller LLMs are still predominantly based on the English language and have rarely been pre-trained on lesser-used, low-resource languages.
Therefore, the model size and multilinguality still present a substantial obstacle to using state-of-the-art LLMs in specific industries, where low latency, high data volumes, autonomy, and resistance to outages are crucial.

In this paper, we analyze possible approaches to multi-label text classification (MTLC) for a specific media monitoring dataset that includes low-resourced languages and a large number of labels.
We propose a novel MTLC method that addresses the challenges of multi-linguality, low latency, model size, and label count based on concepts similar to Retrieval Augmented Generation (RAG) but used in the context of MTLC.
Due to the extremely fast pace of research and development in the field of NLP, we tried to employ a method agnostic to a specific LLM model selection, which wouldn't need constant fine-tuning because of frequent label changes.
This method could also provide a basis for further development of multipurpose systems, including online clustering and summarization for the media monitoring industry.
